\chapter*{How to use this manual}
\addcontentsline{toc}{chapter}{How to use this manual}

This manual is designed for selective, non-linear reading: a reader with a
specific question should be able to turn directly to the relevant chapter without
reading the preceding material. To that end we adopt a consistent internal
structure and several navigation aids.

\paragraph{Chapter structure.} Each chapter opens with a short statement of what
the component does and where it sits in the per-step pipeline
(\cref{ch:arch}). It then develops the governing equations and the discrete
formulae as implemented, accompanied by a schematic that shows which mesh
entities (nodes, element centres, vertical interfaces) supply values and how they
are combined. Each substantive claim about the implementation is followed by a
\textbf{Source} callout naming the file and the function or line range that
realizes it, so that the manual and the code can be read side by side.

\paragraph{Finding a topic.} Three routes are provided. The table of contents
lists every section. \Cref{app:codemap} is a source-code map: it lists, per
subsystem, the authoritative source symbols and the manual section that explains
them, so a reader starting from the code can find the corresponding text. The
index collects symbols, configuration keys, and function names.

\paragraph{Conventions and symbols.} The sign and unit conventions used
throughout (positive-downward depth, the integration-direction factor, the
spherical metric, the distinction between fields stored at nodes, at element
centres, and at vertical interfaces) are stated once in the front matter section
that follows, and the complete symbol list is given in \cref{app:notation}.

\paragraph{Worked configuration.} Where a concrete size clarifies a choice, we refer to the FESOM2
\emph{core2} global mesh; the ready configurations shipped with the engine cover two-dimensional global runs,
a three-dimensional regional run, a structured-grid run and HEALPix runs. The examples illustrate the engine;
they are not the subject of this document.

\paragraph{Setting up a run.} A reader who wants to run an experiment rather than study the numerics can start
from \cref{ch:skills}, which describes how a described experiment becomes seeds, a configuration and a run
command, and how the output is turned into statistics and figures; \cref{app:config} lists every configuration
key.
